\pdfbookmark[1]{Explanations of the plots and numerical readings for the scheduled observations.}{erlaeuterungen-termine}
\begin{center}\large\textit{Explanations of the plots and numerical readings for the scheduled observations.}\par
[\textit{An excerpt.}]\end{center}
\begin{center}\small Results from the Observations of the Magnetic Association.\quad 1836.\quad V.\end{center}

Here ---\quad---\quad---\quad---\quad--- are plots of the observations of magnetic variation made in six scheduled sessions, forty-six curves in all from fourteen different observing sites: Berlin, Breda, Breslau, Catania, Freiberg, Göttingen, The Hague, Leipzig, Milan, Marburg, Messina, Munich, Palermo and Uppsala.
\begin{center}\small
---\enspace---\enspace---\enspace---\enspace---\enspace---\enspace---\enspace---\enspace---\enspace---\enspace---\enspace---\enspace---\enspace---\enspace---
\end{center}
On 28 November 1835 and during the following night, an exceedingly violent sirocco seriously disrupted the observations in Palermo. At one point they even had to be interrupted for an hour and a half:
\SourceLeaf{575}for many sets of readings, only isolated, unreliable determinations could be obtained. There is reason to suppose, therefore, that many of the resulting fluctuations were not genuinely magnetic movements. Even so, we did not want to leave out this curve. Its last part, from the morning of 29 November, when the storm had largely subsided, agrees quite satisfactorily with the records from the northern observing sites.

It may be worth noting here. that, according to all other experience, even the fiercest gales have no effect at all on the magnetic needle, provided that the observing room and the case are well enough sealed to keep out their direct mechanical influence. At the Göttingen magnetic observatory, the needle has very often been observed to remain extremely still, or the variation to proceed very steadily, while a storm of the greatest violence raged outside. But anyone who, on the basis of such experience, was inclined to suppose the very opposite, that storms in the atmosphere acted to paralyze the magnetic forces, would be proved wrong by what happened during the January session of 1836. A very violent storm prevailed in Göttingen and at several other observing sites during that session. When they sent in their results, several observers elsewhere expressed concern that this time the storm might mean that there would be little agreement in the extraordinarily strong movements. Yet during this session, as the plot ---\quad---\quad---\quad--- shows, the curves from the different observing sites agreed so perfectly that one would have to call the agreement astonishing, had so much experience not made it familiar. In several instances here and elsewhere, thunderstorms, even nearby ones, have shown as little discernible effect on the magnetic needle as gales have\expandafter\def\csname HistoricalNoteText575\endcsname{Of course, this does not refer to the case in which atmospheric electricity is conducted to earth by a lead passing through a multiplier surrounding the needle.}\HistoricalNote{575}.

A letter received from Mr.\ von \textsc{Humboldt} in August 1836 reported that from 10---18 August, at Reykjavík in Iceland, an experienced French astronomer, Mr.\ \textsc{Lottin}, would observe magnetic variation with an instrument by \textsc{Gambey}, without interruption, at quarter-hour intervals. The letter expressed the wish that corresponding observations with magnetometers be
\SourceLeaf{576}made on one or several of those days. Accordingly, an extraordinary session was arranged for 17---18 August, and members of our association elsewhere were invited to take part as far as the short notice allowed. This extraordinary session was held in Uppsala, The Hague, Göttingen, Berlin, Leipzig and Munich in exactly the manner introduced for the regular sessions. And while the observations plotted ---\quad---\quad---\quad---\quad--- show some quite interesting movements, we can only regret that the space ---\quad---\quad---\quad--- set aside for the French observations in Iceland had to be left empty, since we were unable to learn anything about the outcome of those French observations.

The September session offers an experience worth mentioning here in some detail, since it confirms in an instructive way the observation made above [p.\ 552 of this volume]. The Marburg observations on that occasion were made in Professor \textsc{Gerling}’s absence, without his personal participation. The record gives only a thoroughly erratic series of numbers for $12^{\mathrm h}5'$, yielding no result at all. At $12^{\mathrm h}10'$ a number suddenly appears that is $30{,}54$ scale divisions higher than the one at $12^{\mathrm h}0'$. ---\quad---\quad---\quad---\quad---\quad---\quad---\quad---\quad---\quad---\quad---\quad--- This led to the conjecture that around $12^{\mathrm h}5'$ a spider had hindered the free movement of the needle by attaching a thread. The conjecture was made still more probable by the fact that, from $12^{\mathrm h}10'$ to the end, the needle’s movements were quite similar to those shown by observations elsewhere, yet proportionally much smaller than one would have had to expect from experience at other sessions. Professor \textsc{Gerling} was therefore asked to inspect the apparatus carefully on his return to Marburg. The result may be added here from a letter from Professor \textsc{Gerling} dated 12 November.

The inspection was made on 5 November. Until then, no one had been back in the observing room since the September session. The position of the needle was determined first and found to be
\begin{center}
\begin{tabular}{rrc l}
$3^{\mathrm h}$&$33'$&.\quad.\quad.\quad.\quad.\quad.&$445{,}63$ scale divisions\\
&35&.\quad.\quad.\quad.\quad.\quad.&$445{,}73$\\
&37&.\quad.\quad.\quad.\quad.\quad.&$445{,}71$
\end{tabular}
\end{center}

\SourceLeaf{577}The needle was then set into moderate oscillation with the aid of the so-called damping rod. This gave a swing duration of 17 seconds, nine seconds less than the previously known swing duration. When the lid of the case was then carefully lifted off, a very small live spider was seen on its underside. A thread, though barely perceptible, also seemed to be hanging from it. A number of small black specks were also found in the case. Under the microscope they proved to be the remains of gnats. Finally, in one corner of the case, there was a full, intact web, so fine that it would hardly have been visible without the reflection of the lights. In view of all the circumstances, one could only suppose that the spider had been living in the case for some time.

After a finger had also been passed around the magnet bar on all sides, new observations of the swing duration gave exactly the old value of 26 seconds again. The position was again found to be considerably lower than before, namely
\begin{center}
\begin{tabular}{rrc l}
$4^{\mathrm h}$&$43'$&.\quad.\quad.\quad.\quad.\quad.&$431{,}15$ scale divisions\\
&45&.\quad.\quad.\quad.\quad.\quad.&$431{,}46$\\
&47&.\quad.\quad.\quad.\quad.\quad.&$431{,}12$
\end{tabular}
\end{center}
Of course, these position readings cannot tell us exactly how much the obstruction that had now been removed had distorted the position, since any change in declination during the interval of more than an hour remained unknown.

This may be a good occasion to mention a second incident of a similar kind. The swing duration of the magnet bar in Breslau, which in March 1836 had been nearly 32 seconds, had increased very gradually from then until November, gaining about $0''4$ in all over that time. This is quite as it should be, since all magnet bars lose some of their strength over time, though in very different proportions, depending on how much the steel is hardened and on other circumstances. Yet from November 1836 to January 1837\expandafter\def\csname HistoricalNoteText577\endcsname{Presumably the swing duration was not determined in the \textit{intervening time}.}\HistoricalNote{577}, on the contrary, there had been a
\SourceLeaf{578}decrease in swing duration of $1''27$. And Professor von \textsc{Boguslawski}, who reported this striking circumstance to me in a letter of 5 March, seemed inclined to attribute it partly to an increase in the intensity of the Earth’s magnetism. To me, however, there seemed no doubt that the cause of this phenomenon had to be sought in the immediate surroundings of the magnet bar, probably in some restriction on its freedom of movement, although the experience at Marburg suggested that a spider’s thread would have had a considerably stronger effect. Mr.\ von \textsc{Boguslawski} also found this conjecture confirmed. On 21 March he replied: ‘You have guessed the cause of the change in swing duration correctly. The case had been accidentally shifted a little to one side, so that the edge of the small hole through which the thread enters from above had come close to the thread, though the two had by no means come into contact. Nevertheless, some fine fibres must have reached the edge, for since the thread has passed through the middle of the hole again, the swing duration has also been nearly the same as before.’

A few remarks may be added here about the movements themselves, as shown for six sessions.

In the three summer sessions ---\quad---\quad---\quad---\quad---\quad---\quad---, the regular daily movement can still be seen through all the large anomalies, in that the curves rise in the afternoon hours and fall during the following morning. In the three winter sessions, however, ---\quad---\quad---\quad---\quad---\quad---\quad---, scarcely anything of that movement can still be discerned. In all our experience it is indeed very common for the irregular movement to overshadow the regular movement or swallow it up altogether. Yet in 1834 and 1835 there were also a few sessions in which no larger anomalies obscured the regular course, though small ones were never absent.

What makes the anomalous movements so remarkable, though, is their extraordinary agreement at different sites, indeed usually at every site, extending as a rule to the smallest nuances, with only the relative magnitudes differing. It would be quite unnecessary to demonstrate this agreement in detail. The plots of our six sessions speak sufficiently for themselves.

\SourceLeaf{579}It cannot yet be our task to decipher this enigmatic hieroglyphic script of nature. For the time being we must confine our efforts to collecting copies of what presents itself and making them ever more reliable, faithful and varied. We may confidently hope that with a wealth of material, decipherment too will come one day. In the meantime we may be allowed to add a few remarks that could contribute to a sounder judgment.

First of all, we must remember that all such anomalies are comparatively slight alterations of or additions to the great magnetic force of the Earth (or, strictly speaking, to its horizontal component), that we must carefully distinguish the anomalies from that force, and that, as things stand, nothing compels us to attribute both to the same or similar causes. One may well consider it probable --- a question we leave entirely open --- that these anomalies could be effects of electric currents or processes of equalization, perhaps far outside the atmosphere. That need not make us abandon the older view that the main force has its seat in the solid part of the Earth itself, or rather that it is the combined effect of all the magnetized parts of the Earth. If, in the view of some naturalists, the interior of the Earth were still liquid, the continuing solidification and consequent thickening of the solid crust would offer the most natural explanation of the secular changes in the magnetic force.

But we would rather leave the unsteady ground of hypotheses and return to the facts. By far the greater number of anomalies are, we find, smaller at the more southerly observing sites and larger at the more northerly ones. For example, the remarkable rise on 30 January 1836 between $9^{\mathrm h}25'$ and $9^{\mathrm h}40'$, reduced to angular units, amounts to 6 in Catania, 12 in Milan, $13\frac12$ in Munich, 16 in Leipzig, 20 in Marburg, 26 in Göttingen and 29 in The Hague, in minutes of arc. Something must indeed be allowed for in this inequality, since at the more northerly points, where the horizontal component of the Earth’s magnetic force itself has a lower intensity than at the more southerly ones, equal disturbing forces must have a greater effect than at the latter. Yet the difference in intensity from The Hague to Catania is only small compared with the observed inequalities. It is therefore certain that
\SourceLeaf{580}the strength of the disturbing force on that occasion was lower the farther south we follow its effect. However little we know about the nature of such disturbing forces, we cannot doubt that each must have its source at a definite place somewhere in space. And just as we must suppose that the source of the force that produced the phenomena just mentioned lay north or northwest of the observing sites (without venturing on a more precise determination from so few data), so in general the farthest northern regions seem to be the chief seat from which the most numerous and greatest effects proceed, insofar as conclusions of this kind can already be based on experience from a region small in comparison with the entire surface of the Earth.

On closer examination of the material available so far, however, one finds considerable differences in the relative magnitudes at different sites of movements that occur in succession, even where the resemblance is otherwise quite unmistakable. Often, for example, of two peaks that follow each other closely, the first is the larger at one site, while the reverse holds at another. We are thus compelled to suppose that on the same day and in the same hour many forces are active at once, that they may be quite independent of one another and have very different seats, and that their effects can mingle in very unequal proportions at different observing sites according to position and distance, or overlap as one force starts to act before another has ceased. It will undoubtedly be very difficult to unravel the complications this introduces into the phenomena at each individual site. Even so, we may confidently hope that these difficulties will not always remain insurmountable once participation in simultaneous observations has spread much more widely. It will be the triumph of science if we one day succeed in bringing order to the varied tangle of phenomena, separating the \textit{individual} forces whose combined result they are, and establishing the seat and magnitude of each.

It is also far from rare to find a small upsurge at one site for which no counterpart can be discerned at the other sites. It would be too bold, however, to declare such a thing at once to be a merely local magnetic
\SourceLeaf{581}effect. With such a large mass of numbers, an error may indeed occur now and then. We have often come across cases in which looking back at the original observations, when we had them in hand, revealed an error in the calculation of the reduction or an obvious slip in writing. In similar cases where we had only an extract of the observations, bringing such a suspicion to the sender’s attention has sometimes had the same result. Since it is impracticable, however, always to discuss all such cases by correspondence first, participants who do not send the original observations themselves may take such points in the curves showing their observations as an occasion to check the originals, as at Leipzig on 26 November 1836 at $6^{\mathrm h}15'$ Göttingen time. Any corrections found in this way will then be announced afterwards. Even then, however, there is still no complete certainty about points of this kind that rest on just one set of observations, even if the original papers fail to establish an error. For even an observer with some practice may occasionally write down the wrong tens digits of the scale readings repeatedly in the same set. Such a conjecture, admittedly a somewhat daring one, would, for example, let us bring the figure of $11.69$ mentioned above down to $6.69$ and so bring it into line with the others.

Even where the fact itself is beyond doubt, such an isolated upsurge still could not be regarded as local in the \textit{strictest} sense. Since the source of every anomaly must have its seat somewhere, the seat of one or another may on occasion be near one of the observing sites itself. If such a force is weak in itself, its effect at the nearest site may be conspicuous precisely because of its proximity, and vanishing (that is, no longer discernible to us) at all the other sites \textit{where observations are made}, precisely because they are already \textit{too} far away. So far at least, therefore, there seems to be no reason at all to suppose that the anomalies differ in anything but quantity. At the same time it follows that in many cases it can be very useful to have two or more observing sites only a moderate distance apart. For example, it would have been very desirable if Augsburg (which now regularly takes part in the scheduled observations) had already observed the September session of
\SourceLeaf{582}1836. Very probably it would then have been possible to judge with greater certainty than we can now the movement at $2^{\mathrm h}10'$, which can be faintly made out at most sites but is strikingly greater in Munich.
\begin{center}\small
---\enspace---\enspace---\enspace---\enspace---\enspace---\enspace---\enspace---\enspace---\enspace---\enspace---\enspace---\enspace---\enspace---\enspace---
\end{center}

